Safety tuning is the main defence that open-weight chat models have against misuse. A worrying possibility is that it works like gradient masking in adversarial robustness: the model learns a surface refusal reflex keyed to trigger words, English, and direct phrasing, which looks robust on familiar inputs but fails under simple distribution shifts. To separate keyword effects from paraphrase and meaning effects, we test 14 open-weight chat models from 0.5B to 32.8B parameters on the same 250 harmful requests under four shifts: keyword removal with a matched re-wording control, translation into nine languages, keyword-\versus-harm decorrelated prompt pools, and role-play or multi-turn framing. We label responses with a validated LLM judge and the \harmbench classifier. Refusal tracks meaning, not vocabulary: its correlation with semantic harm (0.27--0.63) exceeds its correlation with trigger-word count (0.10--0.25) in every model. Removing trigger words lowers refusal by 16.3 percentage points on average, against 2.5 points for a re-wording that keeps them, yet harmful completions rise by only 3.5 points. In sub-3B models, 44--78\% of prompts refused in English are not refused in other languages, but only 6.6--14.8\% receive a reply that carries out the request. The real blind spots are framings: a one-paragraph fiction frame raises compliance from at most 11\% to 51--78\% in four models, and a ``no-rules persona'' raises Llama-3.1-8B to 93\% and Gemma-2-9B to 86\%. These gaps depend on the alignment recipe and do not shrink with scale. Red-teaming should prioritise framing attacks, and multilingual jailbreak rates should count harmful compliance, not non-refusal.
